\documentclass[10pt]{article}
\usepackage[T1]{fontenc}
\usepackage[utf8]{inputenc}
\usepackage[french]{babel}
\usepackage{amsmath,amssymb,mathrsfs,mathpazo}
\usepackage{geometry}
\geometry{a4paper,margin=2cm}
\usepackage{graphicx,array,tabularx,tabularray,colortbl,multirow,multicol,booktabs,enumitem,pifont,ulem,fancybox,diagbox,cancel}
\usepackage{pstricks,pstricks-add,pst-plot,pst-tree,pst-node,pst-text,pst-func,pst-math,pst-bspline,pst-3dplot,pst-eucl,pst-solides3d}
\let\pscalculate\undefined
\usepackage{pst-bezier}
\usepackage{tikz,pgfplots,tkz-tab}
\usetikzlibrary{arrows,arrows.meta,patterns,calc,positioning,trees,decorations.markings,shapes.geometric,shapes.misc,shapes.arrows,intersections,decorations.pathmorphing,decorations.pathreplacing,backgrounds,shadows,babel}
\pgfplotsset{compat=1.18}
\usepgfplotslibrary{fillbetween}
\usepackage[np]{numprint}
\usepackage{listings,makecell,pas-tableur,twoopt}
\let\barre\undefined
\usepackage[pstricks]{tablvar}
\definecolor{darkgreen}{RGB}{0,100,0}
\providecommand{\red}{\color{red}}
\providecommand{\blue}{\color{blue}}
\providecommand{\green}{\color{green}}
\providecommand{\black}{\color{black}}
\providecommand{\white}{\color{white}}
\providecommand{\gray}{\color{gray}}
\providecommand{\cyan}{\color{cyan}}
\providecommand{\magenta}{\color{magenta}}
\providecommand{\yellow}{\color{yellow}}
\providecommand{\violet}{\color{violet}}
\providecommand{\orange}{\color{orange}}
\providecommand{\euro}{\texteuro}
\providecommand{\EUR}{\texteuro}
\providecommand{\ieme}{\textsuperscript{e}}
\providecommand{\degres}{\ensuremath{^\circ}}
\providecommand{\degree}{\ensuremath{^\circ}}
\providecommand{\vv}[1]{\overrightarrow{#1}}
\providecommand{\Coord}[1]{#1}
\providecommand{\decosix}{}
\providecommand{\starredbullet}{\ensuremath{\bullet}}
\pagestyle{empty}
\setlength{\parindent}{0pt}
\usepackage{adjustbox}\DeclareUnicodeCharacter{2236}{:}\DeclareUnicodeCharacter{2190}{\ensuremath{\leftarrow}}
\providecommand{\R}{\mathbb{R}}
\providecommand{\N}{\mathbb{N}}
\providecommand{\Z}{\mathbb{Z}}
\providecommand{\Q}{\mathbb{Q}}
\providecommand{\D}{\mathbb{D}}
\providecommand{\C}{\mathbb{C}}
\providecommand{\e}{{\mathrm{e}}}
\providecommand{\vect}[1]{\overrightarrow{#1}}
\providecommand{\vectt}[1]{\overrightarrow{\mathrm{#1}}}
\providecommand{\vvt}[1]{\overrightarrow{\mathrm{#1}}}
\providecommand{\np}[2][]{{#2~#1}}
\providecommand{\vv}[1]{{\overrightarrow{#1}}}
\renewcommand{\d}{{\mathrm{d}}}
\providecommand{\degres}{{^\circ}}
\providecommand{\degre}{{^\circ}}
\providecommand{\blue}{\color{blue}}
\providecommand{\red}{\color{red}}
\providecommand{\enskip}{\quad}
\providecommand{\llbracket}{\lbrack\!\lbrack}
\providecommand{\rrbracket}{\rbrack\!\rbrack}
\providecommand{\euro}{{\text{€}}}
\providecommand{\no}{{\text{n°}}}
\providecommand{\interval}[2]{{[#1;#2]}}
\providecommand{\intervFO}[2]{{[#1;#2[}}
\providecommand{\dotfill}{\ldots}

\providecommand{\ds}{\displaystyle}
\providecommand{\textding}[1]{#1}
\providecommand{\Oij}{\(\mathrm{O};\vec{\imath},\vec{\jmath}\)}
\providecommand{\Oijk}{\(\mathrm{O};\vec{\imath},\vec{\jmath},\vec{k}\)}
\providecommand{\Ouv}{\(\mathrm{O};\vec{u},\vec{v}\)}
\providecommand{\intervFF}[2]{\ensuremath{[#1;#2]}}
\renewcommand{\euro}{\ifmmode\text{\texteuro}\else\texteuro\fi}
\DeclareUnicodeCharacter{00B0}{\ifmmode^{\circ}\else\textdegree\fi}
\DeclareUnicodeCharacter{20AC}{\ifmmode\text{\texteuro}\else\texteuro\fi}
% Énoncés : Éducation nationale. Transcription et corrigés : APMEP, Denis Vergès et contributeurs. Découpage et classement : CoopMaths, Sébastien Lozano, Éric Elter et contributeurs. Adaptation de la présentation : Maths Explorer.
% CC BY-SA : https://coopmaths.fr/www/about/
% https://coopmaths.fr/alea/static/bac/2026/tex/bac_2026_06_sujet1_metropole_2.tex
\begin{document}
\noindent\textit{APMEP / CoopMaths. CC BY-SA. Présentation adaptée par Maths Explorer.}\par\medskip
Pour chacune des affirmations suivantes, indiquer si elle est vraie ou fausse. Justifier chaque réponse. Une réponse non justifiée ne rapporte aucun point.

	\begin{enumerate}[itemsep=1cm]
		\item L'espace est muni d'un repère orthonormé \Oijk{}.

		On considère :

		\begin{itemize}[label={\textbullet~}]
			\item la droite $(d)$ dont une représentation paramétrique est :
		\end{itemize}
		\[\begin{cases}
				x=t \\
			y=-1,5-t \\
			z=2-2 t
		\end{cases}\text { où } t \in \R; \]
		\begin{itemize}[label={\textbullet~}]
			\item les points A$(3 ; 0 ; 2)$ et B$(2 ; 1 ;-3)$;

			\item le plan $(\mathcal{P})$ d'équation : \quad $-x+y-5 z-0,5=0$.
		\end{itemize}

		\begin{enumerate}[itemsep=5mm]
			\item \textbf{Affirmation 1 :}

			Le plan $(\mathcal{P})$ est orthogonal à la droite (AB), et passe par le milieu du segment [AB].

			\item \textbf{Affirmation 2 :}

			Les droites $(d)$ et (AB) sont sécantes.

			\item On considère le point C de coordonnées $(1,5 ; -3 ; -1)$.

			\textbf{Affirmation 3 :}

			La mesure de l'angle $\widehat{\mathrm{ACB}}$, arrondie à $10^{-1}$, est égale à $70,5$\,\textdegree.
		\end{enumerate}

		\item Titouan et Clotilde participent à un escape game. Ils se retrouvent dans une salle possédant deux portes de sortie, A et B, chacune protégée par un digicode équipé de 8 touches portant des symboles différents :

		\begin{itemize}[label={\textbullet~}]
			\item le digicode de la porte A utilise un code de 3 symboles différents devant être saisis dans l'ordre;
			\item celui de la porte B utilise un code de 4 symboles différents qui peuvent être saisis dans n'importe quel ordre.
		\end{itemize}

		N'étant pas parvenus à obtenir d'indices, ils décident de procéder au hasard à la saisie des codes.

		Clotilde choisit un code pour la porte A, et Titouan choisit un code pour la porte B.

		On considère que les choix des codes sont équiprobables.

		\textbf{Affirmation 4 :}

		Titouan a plus de chances d'ouvrir sa porte que Clotilde.
	\end{enumerate}

\end{document}
